\frontsection{Календарний план проходження практики}

\noindent\textbf{Студент:} Сергійчук Сергій Миколайович, група ІПЗ-61м, спеціальність 121 «Інженерія програмного забезпечення».

\noindent\textbf{База практики:} кафедра інженерії програмного забезпечення та штучного інтелекту НУВГП.

\noindent\textbf{Строк проходження практики:} з 10 серпня 2026 р. по 05 вересня 2026 р.

\vspace{1em}
\begin{spacing}{1.0}
\small\noindent
\renewcommand{\arraystretch}{1.35}
\setlength{\tabcolsep}{6pt}
\begin{xltabular}{\textwidth}{| c | >{\raggedright\arraybackslash}X | >{\centering\arraybackslash}p{2.5cm} | >{\centering\arraybackslash}p{2.3cm} |}
    \hline
    \textbf{№} & \textbf{Назва етапу (виду робіт)} & \textbf{Термін} & \textbf{Відмітка} \\ \hline
    \endfirsthead

    \hline
    \textbf{№} & \textbf{Назва етапу (виду робіт)} & \textbf{Термін} & \textbf{Відмітка} \\ \hline
    \endhead

    1 & Організаційні збори, інструктаж з охорони праці, отримання завдання на практику & 10.08--11.08 & виконано \\ \hline
    2 & Ознайомлення з предметною областю: планування внутрішніх електромереж приміщень, розміщення електрообладнання та склад оздоблювальних матеріалів & 12.08--14.08 & виконано \\ \hline
    3 & Аналіз проблематики планування електромережі та розрахунку матеріалів; формулювання проблеми та мети кваліфікаційної роботи & 17.08--19.08 & виконано \\ \hline
    4 & Огляд наявних засобів (Sweet Home 3D, RoomSketcher, Planner 5D, онлайн-калькулятори гіпсокартону) та літератури із задач розкрою листових матеріалів & 20.08--21.08 & виконано \\ \hline
    5 & Аналіз програмної реалізації вебзастосунку: структура модулів, геометрична модель приміщення, модель даних і формат збереження проєкту & 25.08--26.08 & виконано \\ \hline
    6 & Розроблення специфікації вимог за ISO/IEC/IEEE 29148:2018, побудова моделі варіантів використання & 27.08--28.08 & виконано \\ \hline
    7 & Формалізація архітектури (діаграми C4) та алгоритмів: трасування проводки, розрахунки матеріалів, розкрій гіпсокартону & 31.08--01.09 & виконано \\ \hline
    8 & Верифікація розрахунків: незалежний перерахунок контрольної сцени, вимірювання часу розкрою, порівняння режимів розкладки & 02.09--03.09 & виконано \\ \hline
    9 & Оформлення розділів 1 і 2 кваліфікаційної роботи та списку використаних джерел & 04.09 & виконано \\ \hline
    10 & Підготовка звіту з практики та його захист & 05.09 & виконано \\ \hline
\end{xltabular}
\end{spacing}

\vspace{1.5em}
\signatureline
{Керівник практики від кафедри:}{Жуковський В. В.}
{Студент:}{Сергійчук С. М.}
